\documentclass[../../main.tex]{subfiles}
\graphicspath{{../../image/}{../../figure/}}
\begin{document}

\section{熵产生}\label{section:ch11-熵产生}

鉴于熵耗散在双曲守恒律理论中的中心地位, 确定熵汇的位置十分重要。我们已看到, 在 $BV$ 解范围内熵仅在激波集上产生: 测度\cref{chapter:ch04-cauchy问题}在近似连续点集上消失。然而, 有迹象表明熵浓度与解的特征速度的 $BV$ 结构相联系, 而非与解本身的 $BV$ 结构相联系。确实, 已有证明(见\cref{chapter:ch11-一维标量守恒律}参考文献): 在空间一维标量守恒律情形, 熵仅可能在 $L^\infty$ 解的间断点处产生。为在简单设置中展示这一点, 我们在此证明连续、不必为 $BV$ 的解不产生任何熵。


由\cref{theorem:ch11-11-3-6}, 真正非线性标量守恒律的任意连续解有有界变差。然而, 当通量具有拐点时, 存在无界变差连续解。在某些带凸通量的平衡律设置中也是如此, 源在解上诱导连续性而非有界变差(见\cref{chapter:ch11-一维标量守恒律}参考文献)。


因此考虑守恒律
\begin{align}
\partial_t u(x,t)+\partial_x f(u(x,t))=0,\label{align:ch11-11-13-1}
\end{align}
其中 $f$ 是任意 $C^2$ 函数。设 $\mathscr{U}_-$, $\mathscr{U}_+$ 与 $\mathscr{U}_0$ 分别表示 $u$ 的满足 $f''(u)<0$, $f''(u)>0$ 与 $f''(u)=0$ 的(可能空)集。我们证明如下。

\begin{theorem}\label{theorem:ch11-11-13-1}
假设 $u$ 是 $\mathbb{R}^2$ 某开集 $\mathscr{O}$ 上的连续函数, 在分布意义下满足\eqref{align:ch11-11-13-1}。则
\begin{align}
\partial_t\eta(u(x,t))+\partial_x q(u(x,t))=0\label{align:ch11-11-13-2}
\end{align}
在分布意义下对每个熵-熵流对 $(\eta,q)$ 成立。特别地, 任意连续解及其时间逆转均满足可容许性条件\cref{chapter:ch06-标量守恒律L1理论}。

\end{theorem}

\eqref{align:ch11-11-13-2}的断言是局部的, 故为简化表述, 不妨假设 $u$ 在某条形区域 $(-\infty,\infty)\times[0,s]$ 上定义且连续。上述命题的证明将采用广义特征线理论。与\eqref{align:ch11-11-13-1}连续解 $u$ 关联的特征线是常微分方程 $\dfrac{dx}{dt}=f'(u(x,t))$ 的经典解。在 Lipschitz 连续性缺失时, 标准常微分方程理论保证解的存在性但不保证唯一性。尽管如此, 我们有

\begin{lemma}\label{lemma:ch11-11-13-2}
设 $\xi(\cdot)$ 是与\eqref{align:ch11-11-13-1}连续解 $u$ 关联、定义在 $[0,s]$ 上的任意特征线。则 $u(\xi(t),t)$ 在 $[0,s]$ 上为常数, 因而 $\xi(\cdot)$ 的图像是直线段。特别地, 条形区域 $(-\infty,\infty)\times[0,s]$ 中任意一点通过唯一特征线。
\end{lemma}

\begin{proof}
设 $u(\xi(s),s)=\bar u$。当 $\bar u\in\mathscr{U}_+$ 时, 固定 $r\in[0,s)$ 与 $\varepsilon>0$, 取 $s-r$ 与 $\varepsilon$ 充分小, 使 $u(x,t)\in\mathscr{U}_+$ 在 $\{(x,t):t\in[r,s],\ \xi(t)-\varepsilon\leqslant x\leqslant\xi(t)+\varepsilon\}$ 上。然后对区域 $\{(x,t):r\leqslant t\leqslant s,\ \xi(t)-\varepsilon\leqslant x\leqslant\xi(t)\}$ 与 $\{(x,t):r\leqslant t\leqslant s,\ \xi(t)\leqslant x\leqslant\xi(t)+\varepsilon\}$ 应用守恒律, 得\eqref{align:ch11-11-1-7}与\eqref{align:ch11-11-1-8}。两式右端均非负, 除以 $\varepsilon$ 并令 $\varepsilon\to0$, 推出 $u(\xi(r),r)=u(\xi(s),s)=\bar u$。直接延拓论证进而给出 $u(\xi(t),t)=\bar u$, 对所有 $t\in[0,s]$。

由同样论证, 当 $\bar u\in\mathscr{U}_-$ 时, 即\eqref{align:ch11-11-1-7}与\eqref{align:ch11-11-1-8}右端均非正, 或当 $\bar u\in\mathscr{U}_0$ 的内部时, 即右端均消失, 同样得到结论 $u(\xi(t),t)=\bar u$, $t\in[0,s]$。最后, 由于 $\mathscr{U}_0$ 的边界点不互连, 当 $\bar u$ 位于 $\mathscr{U}_0$ 的边界上时断言 $u(\xi(t),t)=\bar u$, $t\in[0,s]$, 仍成立。证毕。
\end{proof}

由于特征线图像是不相交的直线段, 我们推出

\begin{corollary}\label{corollary:ch11-11-13-3}
对任意 $-\infty<x<y<\infty$ 与 $0<t<s$,
\begin{align}
-\frac{1}{s-t}\leqslant\frac{f'(u(y,t))-f'(u(x,t))}{y-x}\leqslant\frac{1}{t}.\label{align:ch11-11-13-3}
\end{align}
\end{corollary}

下一步是对一类特殊区域建立熵守恒。固定任意光滑熵-熵流对 $(\eta,q)$, $q'(u)=\eta'(u)f'(u)$, 陈述如下。

\begin{lemma}\label{lemma:ch11-11-13-4}
固定 $-\infty<a<b<\infty$ 与 $0<t<\bar t<s$, 设 $u_a=u(a,t)$, $u_b=u(b,t)$, $\bar a=a+(\bar t-t)f'(u_a)$, $\bar b=b+(\bar t-t)f'(u_b)$。则
\begin{align}
\int_{\bar a}^{\bar b}\eta(u(x,\bar t))\,dx-\int_a^b\eta(u(x,t))\,dx\nonumber\\
+(\bar t-t)[q(u_b)-f'(u_b)\eta(u_b)-q(u_a)+f'(u_a)\eta(u_a)]=0.\label{align:ch11-11-13-4}
\end{align}
\end{lemma}

\begin{proof}
取充分小正数 $\varepsilon$, 考虑区间 $[a,b]$ 的任意有限分划, 由点 $a=a_1\leqslant b_1\leqslant a_2\leqslant b_2\leqslant\cdots\leqslant a_n\leqslant a_n=b$ 构成, 满足如下条件。每当 $a_i<b_i$, 开区间 $(a_i,b_i)$ 含于 $\mathscr{U}_+$ 或 $\mathscr{U}_-$。此外,
\begin{align}
\mu\left(\mathscr{U}_0^c\cap\bigcup_{i=1}^{n-1}[b_i,a_{i+1}]\right)<\varepsilon.\label{align:ch11-11-13-5}
\end{align}

对 $i=1,\ldots,n$, 设 $u_i=u(a_i,t)$, $v_i=u(b_i,t)$, 然后 $\bar a_i=a_i+(\bar t-t)f'(u_i)$, $\bar b_i=b_i+(\bar t-t)f'(v_i)$。注意恒等式
\begin{align}
\int_{u_i}^{v_i}f''(w)\eta(w)\,dw=q(u_i)-f'(u_i)\eta(u_i)-q(v_i)+f'(v_i)\eta(v_i),\label{align:ch11-11-13-6}
\end{align}
\begin{align}
\int_{v_i}^{u_{i+1}}f''(w)\eta(w)\,dw=q(v_i)-f'(v_i)\eta(v_i)-q(u_{i+1})+f'(u_{i+1})\eta(u_{i+1}).\label{align:ch11-11-13-7}
\end{align}
由中值定理,
\begin{align}
\int_{u_i}^{v_i}f''(w)\eta(w)\,dw=\eta(w_i)[f'(v_i)-f'(u_i)]\nonumber\\
=\frac{1}{\bar t-t}\eta(w_i)[(\bar b_i-\bar a_i)-(b_i-a_i)],\label{align:ch11-11-13-8}
\end{align}
其中 $w_i$ 位于 $u_i$ 与 $v_i$ 之间。此外, 取 $w_i^*=\frac12(v_i+u_{i+1})$,
\begin{align}
\eta(w_i^*)[f'(u_{i+1})-f'(v_i)]=\frac{1}{\bar t-t}\eta(w_i^*)[(\bar a_{i+1}-\bar v_i)-(a_{i+1}-v_i)].\label{align:ch11-11-13-9}
\end{align}

由\eqref{align:ch11-11-13-5}与\eqref{align:ch11-11-13-3},
\begin{align}
\sum_{i=1}^{n-1}TV_{[b_i,a_{i+1}]}f'(u(\cdot,t))\leqslant c\varepsilon.\label{align:ch11-11-13-10}
\end{align}
特别地,
\begin{align}
\sum_{i=1}^{n-1}|f'(u_{i+1})-f'(v_i)|\leqslant c\varepsilon.\label{align:ch11-11-13-11}
\end{align}

此外, 由 Stieltjes 积分
\begin{align}
\int_{v_i}^{u_{i+1}}f''(w)\eta(w)\,dw=\int_{b_i}^{a_{i+1}}\eta(u(\cdot,t))\,df'(u(\cdot,t))\label{align:ch11-11-13-12}
\end{align}
以及\eqref{align:ch11-11-13-7}、\eqref{align:ch11-11-13-10}蕴含
\begin{align}
\sum_{i=1}^{n-1}|q(v_i)-f'(v_i)\eta(v_i)-q(u_{i+1})+f'(u_{i+1})\eta(u_{i+1})|\leqslant c\varepsilon.\label{align:ch11-11-13-13}
\end{align}

现合并\eqref{align:ch11-11-13-8}、\eqref{align:ch11-11-13-9}、\eqref{align:ch11-11-13-11}与\eqref{align:ch11-11-13-13}, 得
\begin{align}
\biggl|\sum_{i=1}^n\eta(w_i)(\bar b_i-\bar a_i)+\sum_{i=1}^{n-1}\eta(w_i^*)(\bar a_{i+1}-\bar b_i)\nonumber\\
-\sum_{i=1}^n\eta(w_i)(b_i-a_i)-\sum_{i=1}^{n-1}\eta(w_i^*)(a_{i+1}-b_i)\nonumber\\
+(\bar t-t)[q(u_b)-f'(u_b)\eta(u_b)-q(u_a)+f'(u_a)\eta(u_a)]\biggr|\leqslant c\varepsilon,\label{align:ch11-11-13-14}
\end{align}
其中 $c$ 不依赖于 $[a,b]$ 的所选分划或 $\varepsilon$。当 $[a,b]$ 的分划加细时, \eqref{align:ch11-11-13-14}左端前两个和收敛到 $\eta(u(\cdot,\bar t))$ 在 $[\bar a,\bar b]$ 上的积分, 后两个和收敛到 $\eta(u(\cdot,t))$ 在 $[a,b]$ 上的积分。由于 $\varepsilon$ 任意小, 得\eqref{align:ch11-11-13-4}。证毕。
\end{proof}

现固定某 $\tau\in(0,s)$, 考虑坐标变换 $(x,t)\mapsto(z,t)$, 定义为
\begin{align}
z=x-(t-\tau)f'(u(x,\tau)).\label{align:ch11-11-13-15}
\end{align}
由\cref{lemma:ch11-11-13-2}与\cref{corollary:ch11-11-13-3}, 上述变换是双 Lipschitz 同胚, 逆为
\begin{align}
x=z+(t-\tau)f'(u(z,\tau)).\label{align:ch11-11-13-16}
\end{align}

注意在新坐标系中特征线化为 $z=$ 常数, $u$ 成为 $z$ 的函数: $u(x(z,t),t)=\hat u(z)$。因而自然地将 $(x,t)$ 与 $(z,t)$ 分别解释为由\eqref{align:ch11-11-13-1}生成流的"Eulerian"与"Lagrangian"表示。

\textbf{\cref{theorem:ch11-11-13-1}的证明。} 设
\begin{align}
\theta(z,t)=\eta(\hat u(z))\partial_z x(z,t),\label{align:ch11-11-13-17}\\
p(z)=q(\hat u(z))-f'(\hat u(z))\eta(\hat u(z)),\label{align:ch11-11-13-18}
\end{align}
注意\eqref{align:ch11-11-13-4}蕴含
\begin{align}
\int_\alpha^\beta[\theta(z,\bar t)-\theta(z,t)]\,dz+(\bar t-t)[p(\beta)-p(\alpha)]=0,\label{align:ch11-11-13-19}
\end{align}
对所有 $-\infty<\alpha<\beta<\infty$ 与 $0<t<\bar t<s$。于是 $p$ 在 $(-\infty,\infty)$ 上 Lipschitz, 且
\begin{align}
\frac{\theta(z,\bar t)-\theta(z,t)}{\bar t-t}+\frac{d}{dz}p(z)=0\label{align:ch11-11-13-20}
\end{align}
对几乎所有 $z$ 成立, 亦在分布意义下对任意 $0<t<\bar t<s$ 成立。进而
\begin{align}
\partial_t\theta(z,t)+\partial_z p(z)=0,\label{align:ch11-11-13-21}
\end{align}
在 $(-\infty,\infty)\times(0,s)$ 上分布意义下。

固定任意 $C^\infty$ 试验函数 $\phi(x,t)$, 紧支撑于 $(-\infty,\infty)\times(0,s)$, 设 $\psi(z,t)=\phi(x(z,t),t)$。于是 $\psi$ 是 Lipschitz 函数, 紧支撑于 $(-\infty,\infty)\times(0,s)$, 导数 $\partial_z\psi=\partial_x\phi\partial_z x$, $\partial_t\psi=\partial_t\phi+\partial_x\phi\partial_t x$。考虑到 $\partial_t x=f'(u)$, 并回忆\eqref{align:ch11-11-13-17}与\eqref{align:ch11-11-13-18},
\begin{align}
\theta\partial_t\psi+p\partial_z\psi=\partial_z x[\eta(u)\partial_t\phi+q(u)\partial_x\phi].\label{align:ch11-11-13-22}
\end{align}
因此, 由\eqref{align:ch11-11-13-21}及 $\partial_x z=[\partial_z x]^{-1}$,
\begin{align}
\int_0^s\int_{-\infty}^\infty[\eta(u)\partial_t\phi+q(u)\partial_x\phi]\,dxdt
=\int_0^s\int_{-\infty}^\infty[\theta\partial_t\psi+p\partial_z\psi]\,dzdt=0,\label{align:ch11-11-13-23}
\end{align}
这建立了\eqref{align:ch11-11-13-2}。证毕。

\end{document}
